%% ===================================================================
\section{DRTO 모델 구조}
\label{sec:drto-model}
%% ===================================================================

앞서 도출한 이론적 배경과 선행 연구의 한계를 토대로, 본 절부터는 동적
복구목표시간($\DRTO$) 프레임워크의 수학적 구조와 설계 원리를 체계적으로 기술한다.
이후 계속하여 구체적으로 (i) 시그모이드 정식화(2.7), (ii) 4단계 조건 체계 C0--C3(2.8),
(iii) 수학적 성질(2.9), (iv) 곡률 매개변수 $\kappa^{*}=1.2$ 선정 근거(2.10),
(v) \jrev{10개 연구가설} H1--\jrev{H10}(2.11)을 차례로 다룬다.

%% -------------------------------------------------------------------
\subsection{\chg{개별 시그모이드 바운딩 핵심 수식}}
\label{subsec:core-equation}
%% -------------------------------------------------------------------

$\DRTO$ 모델의 핵심 과제는 관측성과 불확실성이라는 두 입력 차원의 효과를
물리적으로 타당한 범위 내에서 결합하는 것이다. 복구목표시간은 음수가 될 수 없으며,
조직이 허용할 수 있는 최대허용중단기간(MTPD)을 초과해서도 안 된다.
이러한 물리적 제약 조건을 만족하면서도 입력 변수의 변화에 연속적이고
미분 가능한 응답을 제공하는 모델을 구축하는 것이 설계의 핵심 요구사항이다.

본 연구에서는 \textbf{개별 시그모이드 바운딩(individual sigmoid bounding)} 모델을 제안한다.
선행연구\citep{lee2026drto}에서 수립한 관측성 기반 모델(C0$\to$C1)과
최적 곡률 $\kappa^{*} = 1.2$를 기반으로,
본 연구에서는 불확실성 채널(C2, C3)을 추가하여 다채널 구조로 확장한다.
관측성과 불확실성의 효과를 각각 고유한 물리적 범위(span)에서
독립적으로 바운딩하여, 물리적 경계 보장, 연속 미분 가능성,
이론적 유도 가능성이라는 세 요구사항을 동시에 충족한다.

\chg{모델의 기저 함수는} 수정 시그모이드(modified sigmoid)로서,
표준 로지스틱 함수를 $(-1, 1)$ 구간으로 사상한다.

\begin{equation}
  S(z) = \frac{2}{1 + e^{-z}} - 1, \quad S: \mathbb{R} \to (-1,\; 1)
  \label{eq:sigmoid}
\end{equation}

\chg{관측성의 효과는} 시그모이드 함수를 통해 $(-R_0, 0]$ 범위로 바운딩된다\chg{. 이때 $0$은
달성 가능한 값이지만 $-R_0$에는 한없이 접근할 뿐 도달하지는 못한다}.
관측성은 $\DRTO$를 $R_0$에서 최대 $0$까지 감소시킬 수 있으므로,
물리적 하향 범위(span)는 $R_0$이다.
\chg{이와 대칭적으로 불확실성의 효과는} 시그모이드 함수를 통해 $[0, R_{\max} - R_0)$ 범위로 바운딩된다\chg{. $0$은
달성 가능하나 $R_{\max} - R_0$에는 도달하지 못하며, $R_0 = 6.0$h와 $R_{\max} = 24.0$h를
가정하면 그 범위는 $[0, 18.0)$h이다}.
불확실성은 $\DRTO$를 $R_0$에서 최대 $R_{\max}$(MTPD)까지 증가시킬 수 있으므로,
물리적 상향 범위(span)는 $R_{\max} - R_0$이다.

\medskip
두 채널의 바운딩 수식은 다음과 같이 대칭적 구조를 갖는다.
\begin{align}
  E_{O,\text{bnd}} &= -R_0 \cdot S(z_O),
    &\quad \text{span} &= R_0
    \label{eq:E_O_bnd} \\[4pt]
  E_{U,\text{bnd}} &= \;(R_{\max} - R_0) \cdot S(z_U),
    &\quad \text{span} &= R_{\max} - R_0
    \label{eq:E_U_bnd}
\end{align}
시그모이드 입력 인자 $z_O$, $z_U$는 각각
\begin{align}
  z_O &= \kappa \cdot k_O \cdot O_{\text{raw}}
    \label{eq:z_O} \\[4pt]
  z_U &= \frac{\kappa \cdot k_U \cdot R_0 \cdot U_{\text{raw}}}{R_{\max} - R_0},
    \quad k_U = 1 - k_O
    \label{eq:z_U}
\end{align}
이다. $z_O = \kappa \cdot k_O \cdot O_{\text{raw}} \geq 0$은
입력 변수의 정의역에 의해 항상 성립하며,
$U_{\text{raw}} \geq 0$일 때 $z_U \geq 0$이므로
$E_{O,\text{bnd}} \in (-R_0, 0]$,\;
$E_{U,\text{bnd}} \in [0, R_{\max} - R_0)$이다.
$N(3,1)$에서의 음수 $U_{\text{raw}}$ 처리에 대해서는
\ref{subsec:c2}항에서 후술한다.

\chg{두 입력 인자 $z_O$, $z_U$는 세 가지 원칙에 따라 설계된다. 첫째는 무차원
정규화이다. 관측성 점수와 불확실성 원시값은 단위와 척도가 서로 다르므로, 각각
가중치를 곱한 뒤 불확실성은 물리적 상향 범위 $R_{\max} - R_0$로 나누어 두 입력을
모두 무차원량으로 변환한다. 이로써 시그모이드가 비교 가능한 척도 위에서
작동한다. 둘째는 물리적 범위 분리이다. 관측성은 복구 시간을 기준선 아래로
($\mathrm{span} = R_0$), 불확실성은 기준선 위로($\mathrm{span} = R_{\max} - R_0$)
작용하도록 두 채널의 방향과 범위를 분리한다. 셋째는 상보적 가중이다. 관측성
가중치와 불확실성 가중치가 $k_O + k_U = 1$을 이루어, 한 채널의 비중이 커지면
다른 채널의 비중이 줄어드는 상보 관계를 갖는다.}

두 입력 인자는 표면적 형태는 다르나 동일한 구조 $z = \kappa \cdot k \cdot (R_0/\mathrm{span})
\cdot X_{\text{raw}}$를 공유한다. 관측성 채널은 $\mathrm{span} = R_0$이므로 정규화 인자가
$R_0/R_0 = 1$이 되어 수식에서 생략되고, 불확실성 채널은 $\mathrm{span} = R_{\max} - R_0$이므로
정규화 인자가 $R_0/(R_{\max}-R_0) = 1/3$로 나타난다. 두 원시 입력은 모두 무차원이며,
시간 단위(h)는 바운딩 수식의 span에 의해 부여된다.
\p{[}표~\ref{tab:zo-zu-symmetry}\p{]}은 두 채널의 항목별 대칭 구조를 정리한 것이다.

\begin{table}[htbp]
\centering
\caption{관측성 채널($z_O$)과 불확실성 채널($z_U$)의 대칭 구조 비교}
\label{tab:zo-zu-symmetry}
\small
\begin{tabular}{l c c}
\toprule
\textbf{구분} & \textbf{관측성 ($z_O$)} & \textbf{불확실성 ($z_U$)} \\
\midrule
원시 입력     & $O_{\text{raw}} \in [0,1]$ (무차원) & $U_{\text{raw}} \geq 0$ (무차원) \\
정규화 인자   & $R_0 / R_0 = 1$              & $R_0 / (R_{\max} - R_0) = 1/3$ \\
인자 수식     & $\kappa \cdot k_O \cdot (R_0/R_0) \cdot O_{\text{raw}}$
              & $\kappa \cdot k_U \cdot (R_0/(R_{\max}\!-\!R_0)) \cdot U_{\text{raw}}$ \\
가중치        & $k_O$                               & $k_U = 1 - k_O$ \\
span (바깥)   & $R_0$                               & $R_{\max} - R_0$ \\
효과 방향     & 하향 ($\DRTO$ 감소)                 & 상향 ($\DRTO$ 증가) \\
\bottomrule
\end{tabular}
\end{table}

\chg{최종 $\DRTO$는} 기준선 $R_0$에 개별 바운딩된 두 효과를 가산하여 산출한다.

\begin{equation}
  \DRTO = R_0 + E_{O,\text{bnd}} + E_{U,\text{bnd}}
  \label{eq:drto}
\end{equation}

두 바운딩 효과의 범위가 물리적으로 분리되어 있으므로,
$\DRTO$는 항상 다음의 경계 내에 위치한다.

\begin{equation}
  E_{O,\text{bnd}} \in (-R_0, 0], \quad
  E_{U,\text{bnd}} \in [0, R_{\max} - R_0)
  \quad \therefore\;\; \DRTO \in (0,\; R_{\max})
  \label{eq:bounds}
\end{equation}

이 경계는 $\kappa$ 값에 무관하게 항상 성립하며,
$R_{\max} = 4R_0 = 24.0$h로 설정할 경우 $\DRTO \in (0,\; 24.0)$h가 보장된다.
특히 C1 조건($E_{U,\text{bnd}} = 0$)에서는 $\DRTO \in (0, R_0]$이 되어
$R_{\max}$가 불필요하다. 이는 관측성만의 효과를 불확실성 상한과
독립적으로 분석할 수 있게 하는 개별 바운딩의 핵심 장점이다.


\chg{조건 간 비교를 위해} 무차원 정규화 비율 $\eta$를 정의한다.

\begin{equation}
  \eta = \frac{\DRTO}{R_0}
  \label{eq:eta}
\end{equation}

$\eta = 1$은 정적 기준선과 동일함을, $\eta < 1$은 관측성에 의한 복구 시간 단축을,
$\eta > 1$은 불확실성에 의한 복구 시간 연장을 의미한다.
$\DRTO \in (0, R_{\max})$이므로 $\eta$의 이론적 범위는
$(0,\; R_{\max}/R_0)$이며, $R_{\max} = 4R_0$ 가정 하 $(0,\; 4)$이다.


%% -------------------------------------------------------------------
\subsection{모델 구조 도식}
\label{subsec:model-diagram}
%% -------------------------------------------------------------------

\p{[}그림~\ref{fig:model-structure}\p{]}는 $\DRTO$ 개별 시그모이드 바운딩 모델의
전체 구조를 도식화한 것이다. 두 입력 차원(관측성, 불확실성)이
각각 독립적인 시그모이드 바운딩 경로를 통과한 후 기준선 $R_0$에
가산되어 최종 $\DRTO$를 산출하는 구조를 보여준다.
\chg{그림에서 관측성 경로는 상단에 파란색으로, 불확실성 경로는 하단에 주황색으로 표시되며, 두 경로는 각각 독립적인 span에서 시그모이드 바운딩을 거쳐 가산 노드에서 기준선 $R_0$와 합산된다. 곡률 매개변수 $\kappa = 1.2$는 두 경로에 공통으로 적용된다.}

\begin{figure}[htbp]
\centering
\resizebox{0.95\textwidth}{!}{%
\begin{tikzpicture}[
  inputbox/.style={draw, rounded corners=5pt, minimum width=26mm,
    minimum height=13mm, align=center, font=\small, line width=0.6pt},
  procbox/.style={draw, rounded corners=4pt, minimum width=36mm,
    minimum height=12mm, align=center, font=\small, line width=0.6pt},
  outbox/.style={draw, rounded corners=5pt, minimum width=40mm,
    minimum height=16mm, align=center, font=\small\bfseries,
    line width=0.8pt},
  sumbox/.style={draw, circle, minimum size=10mm, align=center,
    font=\normalsize\bfseries, line width=0.8pt},
  arr/.style={-{Stealth[length=2.5mm]}, thick},
  lbl/.style={font=\scriptsize, fill=white, inner sep=1.5pt}
]

% ── 관측성 경로 (상단, y=3.0) ──
\node[inputbox, fill=blue!10] (ko) at (0, 4.0) {\textbf{가중치} $k_O$\\$\in [0,1]$};
\node[inputbox, fill=blue!10] (oraw) at (0, 2.0) {\textbf{원시 관측성} $O_{\text{raw}}$\\$\in [0,1]$};
\node[procbox, fill=blue!5] (zo) at (5.0, 3.0) {$z_O = \kappa \cdot k_O \cdot O_{\text{raw}}$};
\node[procbox, fill=blue!12] (sobnd) at (10.5, 3.0)
  {$E_{O,\text{bnd}} = -R_0 \cdot S(z_O)$\\{\scriptsize span $= R_0$}};

% ── 기준선 (중앙, y=0) ──
\node[procbox, fill=gray!15] (r0) at (5.0, 0) {기준선 $R_0 = 6.0$h};

% ── 불확실성 경로 (하단, y=-3.0) ──
\node[inputbox, fill=orange!10] (ku) at (0, -1.5) {\textbf{가중치} $k_U$\\$= 1 - k_O$};
\node[inputbox, fill=orange!10] (unorm) at (0, -4.0) {\textbf{원시 불확실성} $U_{\text{raw}}$\\$\geq 0$};
\node[procbox, fill=orange!5] (zu) at (5.0, -3.0)
  {$z_U = \kappa \cdot k_U \cdot \frac{R_0 \cdot U_{\text{raw}}}{R_{\max} - R_0}$};
\node[procbox, fill=orange!12] (subnd) at (10.5, -3.0)
  {$E_{U,\text{bnd}} = (R_{\max}\!-\!R_0) \cdot S(z_U)$\\{\scriptsize span $= R_{\max} - R_0$}};

% ── 가산 노드 ──
\node[sumbox, fill=green!8] (sum) at (14.5, 0) {$+$};

% ── 출력 ──
\node[outbox, fill=green!8] (drto) at (19.0, 0)
  {$\DRTO = R_0 + E_{O} + E_{U}$\\$\eta = \DRTO / R_0 \in (0,\; 4)$};

% ── 화살표: 관측성 경로 ──
\draw[arr, blue!60] (ko.east) -- (zo.150);
\draw[arr, blue!60] (oraw.east) -- (zo.210);
\draw[arr, blue!70] (zo.east) -- (sobnd.west);

% ── 화살표: 불확실성 경로 ──
\draw[arr, orange!60] (ku.east) -- (zu.150);
\draw[arr, orange!60] (unorm.east) -- (zu.210);
\draw[arr, orange!70] (zu.east) -- (subnd.west);

% ── 화살표: 3경로 → 가산 노드 (진입점 분리, 겹침 방지) ──
\draw[arr, blue!70]   (sobnd.east) -- ++(0.8,0) |- (sum.130);
\draw[arr, black]  (r0.east) -- ++(6.5,0) -- (sum.west);
\draw[arr, orange!70] (subnd.east) -- ++(0.8,0) |- (sum.230);

% ── 화살표: 가산 노드 → 출력 ──
\draw[arr, black] (sum.east) -- (drto.west);

% ── 매개변수 κ (중앙, 두 경로 사이) ──
\node[draw, rounded corners=3pt, fill=red!5, inner sep=3pt,
  font=\scriptsize, text=black] at (5.0, 0.9)
  {$\kappa = 1.2$};

% ── κ → 두 z 노드 연결 (점선) ──
\draw[densely dashed, red!40, thin, -{Stealth[length=1.5mm]}]
  (5.0, 0.7) -- (zo.south);
\draw[densely dashed, red!40, thin, -{Stealth[length=1.5mm]}]
  (5.0, -0.7) -- (5.0, -2.2);

% ── 범위 라벨 ──
\node[font=\scriptsize, text=black] at (10.5, 4.5)
  {$E_{O,\text{bnd}} \in (-R_0,\; 0]$};
\node[font=\scriptsize, text=black] at (10.5, -4.5)
  {$E_{U,\text{bnd}} \in [0,\; R_{\max}-R_0)$};

% ── 물리적 경계 ──
\node[draw, dashed, rounded corners=3pt, inner sep=5pt,
  font=\scriptsize, text=black] at (19.0, -2.5)
  {물리적 경계: $0 < \DRTO < R_{\max} = 24.0$h};

\end{tikzpicture}%
}% end resizebox
\caption{$\DRTO$ 개별 시그모이드 바운딩 모델 구조}
\label{fig:model-structure}
\end{figure}


%% -------------------------------------------------------------------
\subsection{변수 및 매개변수 정의}
\label{subsec:parameters}
%% -------------------------------------------------------------------

\p{[}표~\ref{tab:parameters}\p{]}은 $\DRTO$ 모델의 전체 변수와 매개변수를 정리한 것이다.
각 매개변수의 설정 근거는 연구자의 논문 \citet{lee2026drto}, 국제 표준, 또는 수학적 제약 조건에
기반한다.

\begin{table}[htbp]
\centering
\caption{DRTO 모델 변수 및 매개변수 정의}
\label{tab:parameters}
\small
\begin{tabularx}{\textwidth}{l l c c X}
\toprule
\textbf{기호} & \textbf{명칭} & \textbf{값/범위} & \textbf{단위} & \textbf{설정 근거} \\
\midrule
$R_0$ & 기준선 RTO & 6.0 & h & \nrev{ISO 22301:2019} BIA 기반 표준 설정값 \\
$R_{\max}$ & DRTO 상한 (MTPD) & 24.0 & h & $R_{\max} = 4 \cdot R_0$ (MTPD 기반 상한) \\
$\kappa$ & 곡률 매개변수 & 1.2 & --- & 다기준 최적화($\kappa^{*}$)~\citep{lee2026drto} \\
$S(z)$ & 수정 시그모이드 & $(-1, 1)$ & --- & $S(z) = 2/(1+e^{-z})-1$ \\
$k_O$ & 관측성 가중치 & $[0, 1]$ & --- & $\mathrm{Uniform}(0, 1)$으로 시뮬레이션 \\
$O_{\text{raw}}$ & 관측성 원시값 & $[0, 1]$ & --- & $\mathrm{Uniform}(0, 1)$으로 시뮬레이션 \\
$k_U$ & 불확실성 가중치 & $[0, 1]$ & --- & $k_U = 1 - k_O$ (상보적) \\
$z_O$ & 관측성 인자 & $\geq 0$ & --- & $\kappa \cdot k_O \cdot O_{\text{raw}}$ \\
$z_U$ & 불확실성 인자 & $\geq 0$ & --- & $\kappa \cdot k_U \cdot R_0 \cdot U_{\text{raw}} / (R_{\max} - R_0)$ \\
$U_{\text{raw}}^{(G)}$ & Gaussian 불확실성 & $N(3, 1)$ & --- & 평균 $\mu=3$, 표준편차 $\sigma=1$ \\
$U_{\text{raw}}^{(P)}$ & Pareto 불확실성 & $6 \cdot \mathrm{Pareto}(3)$ & --- & $\alpha=3$, $x_m=1$, 스케일링($\times 6$) \\
$E_{O,\text{bnd}}$ & 관측성 바운딩 조정량 & $(-R_0, 0]$ & h & span $= R_0 = 6.0$h \\
$E_{U,\text{bnd}}$ & 불확실성 바운딩 조정량 & $[0, R_{\max}\!-\!R_0)$ & h & span $= R_{\max} - R_0 = 18.0$h \\
$\eta$ & 효율비 & $(0, 4)$ & --- & $\DRTO / R_0$ \\
\bottomrule
\end{tabularx}
\end{table}


%% -------------------------------------------------------------------
\subsection{물리적 경계 보장 증명}
\label{subsec:boundary-proof}
%% -------------------------------------------------------------------

\phantomsection\label{prop:boundary}%
임의의 $\kappa > 0$, $k_O \in [0, 1]$, $O_{\text{raw}} \in [0, 1]$,
$z_U \geq 0$에 대해 $0 < \DRTO < R_{\max}$가 성립한다.(명제 2.1)

이 명제의 증명은 다음과 같다.
$z_O = \kappa \cdot k_O \cdot O_{\text{raw}} \geq 0$이므로
$S(z_O) \in [0, 1)$이다. 따라서
$E_{O,\text{bnd}} = -R_0 \cdot S(z_O) \in (-R_0, 0]$이다.
$z_U \geq 0$이므로
$S(z_U) \in [0, 1)$이다. 따라서
$E_{U,\text{bnd}} = (R_{\max} - R_0) \cdot S(z_U) \in [0, R_{\max} - R_0)$이다.
이 두 부등식을 합산하면 $\DRTO = R_0 + E_{O,\text{bnd}} + E_{U,\text{bnd}} \in (0,\; R_{\max})$이므로
$0 < \DRTO < R_{\max}$가 $\kappa$에 무관하게 성립한다. (명제 증명 완료)

위 증명은 $z_U \geq 0$ (즉, $U_{\text{raw}} \geq 0$) 전제 하에 성립한다.
C2 조건에서 $N(3,1)$의 음수 표본($0.135\%$)은 $z_U < 0$을 유발하여
$E_{U,\text{bnd}} < 0$이 되나,
$|E_{O,\text{bnd}}| < R_0$이고 $|E_{U,\text{bnd}}| < R_{\max} - R_0$이므로
$\DRTO > R_0 - R_0 - (R_{\max} - R_0) = 2R_0 - R_{\max}$이다.
$R_{\max} = 4R_0$일 때 이론적 하한은 $-2R_0$이나,
실제 시뮬레이션에서 $\DRTO$ 최솟값은 $2.861$h로 $0$을 충분히 상회한다.
즉 이론 하한이 음수임에도 음수 표본 출현 확률이 $0.135\%$로 작아 실질적 영향은 무시 가능하다.

\medskip
\chg{\noindent 이 경계 보장은 세 가지 실무적 의의를 갖는다. 첫째, 양수
보장($\DRTO > 0$)으로 $\DRTO$가 음수가 되는 물리적으로 무의미한 상황을
원천적으로 방지한다. 둘째, 상한 보장($\DRTO < R_{\max}$)으로
\nrev{ISO 22301:2019}의 MTPD 제약을 수학적으로 충족한다.
셋째, 클리핑(hard bounding)과 달리 경계 근처에서 미분 가능하여 매개변수 변화에
연속적으로 응답한다.}


%% ===================================================================
\section{4개 실험 조건 (C0--C3)}
\label{sec:conditions}
%% ===================================================================

$\DRTO$ 모델의 검증을 위해 관측성과 불확실성의 활성화 조합에 따라
네 개의 실험 조건(C0--C3)을 체계적으로 정의한다.
이 중 C0(기준선)과 C1(관측성 활성화)은 \citet{lee2026drto}에서
\chg{수립되고 검증된} 조건이며, C2(가우시안 불확실성)와 C3(파레토 불확실성)는
본 연구에서 새롭게 도입한 조건이다.
조건 설계의 핵심 원리는 \textbf{점진적 활성화(progressive activation)}로서,
C0에서 C3까지 한 번에 하나의 요소를 추가하여 각 요소의 개별 효과와
결합 효과를 분리 분석할 수 있도록 한다.

\p{[}표~\ref{tab:conditions}\p{]}는 각 조건의 구성과 특성을 비교한 것이다.

\begin{table}[htbp]
\centering
\caption{4개 실험 조건(C0--C3) 정의 및 시뮬레이션 결과 요약}
\label{tab:conditions}
\small
\begin{tabularx}{\textwidth}{c l c c l c X}
\toprule
\textbf{조건} & \textbf{명칭} & \textbf{관측성} & \textbf{불확실성}
  & \textbf{$U_{\text{raw}}$ 분포} & \textbf{평균 $\eta$} & \textbf{수식} \\
\midrule
\textbf{C0} & 정적 기준선 & OFF & OFF & --- & 1.000 &
  $\DRTO = R_0$ \\
\textbf{C1} & 관측성 전용 & ON & OFF & --- & 0.853 &
  $\DRTO = R_0 + E_{O,\text{bnd}}$ \\
\textbf{C2} & Gaussian 불확실성 & ON & ON & $N(3, 1)$ & 1.682 &
  $\DRTO = R_0 + E_{O,\text{bnd}} + E_{U,\text{bnd}}^{(G)}$ \\
\textbf{C3} & Pareto 불확실성 & ON & ON & $6 \cdot \mathrm{Pa}(3)$ & 1.514 &
  $\DRTO = R_0 + E_{O,\text{bnd}} + E_{U,\text{bnd}}^{(P)}$ \\
\bottomrule
\end{tabularx}
\end{table}


\p{[}그림~\ref{fig:conditions-progression}\p{]}은 C0에서 C3까지의 조건 진행과
각 단계에서 활성화되는 요소를 도식화한 것이다.
\chg{정적 조건 C0에서 관측성을 도입한 C1로 이어지는 공통 경로를 거친 뒤, 가우시안 불확실성의 C2와 파레토 불확실성의 C3가 C1에서 병렬로 분기하는 구조이다.}

\begin{figure}[htbp]
\centering
\resizebox{0.95\textwidth}{!}{%
\begin{tikzpicture}[
  cond/.style={draw, rounded corners=6pt, minimum width=30mm,
    minimum height=20mm, align=center, font=\small, line width=0.7pt},
  arr/.style={-{Stealth[length=3mm]}, thick},
  arrlbl/.style={font=\scriptsize, fill=white, inner sep=2pt,
    rounded corners=1pt},
  comp/.style={font=\scriptsize, text=black, align=center},
  efflbl/.style={font=\scriptsize, align=center, text=black}
]

% 조건 노드
\node[cond, fill=gray!15] (c0) at (0, 0)
  {\textbf{C0}\\정적 기준선\\$\eta = 1.000$\\$\DRTO = 6.0$h};
\node[cond, fill=blue!12] (c1) at (5.5, 0)
  {\textbf{C1}\\관측 전용\\$\bar{\eta} = 0.853$\\$\overline{\DRTO} = 5.12$h};
\node[cond, fill=green!12] (c2) at (11, 0)
  {\textbf{C2}\\Gaussian 불확실성\\$\bar{\eta} = 1.682$\\$\overline{\DRTO} = 10.09$h};
\node[cond, fill=red!12] (c3) at (16.5, 0)
  {\textbf{C3}\\Pareto 불확실성\\$\bar{\eta} = 1.514$\\$\overline{\DRTO} = 9.09$h};

% 화살표 및 라벨 (상단: 전이 조건, 중앙: Cohen's d)
\draw[arr] (c0) -- (c1);
\node[arrlbl, above=3pt] at ($(c0)!0.5!(c1)$) {$+O_{\text{raw}},\;k_O$};
\node[efflbl] at ($(c0)!0.5!(c1) + (0, 1.3)$)
  {Cohen's $d = -1.180$\\(C1 vs C0)};

\draw[arr] (c1) -- (c2);
\node[arrlbl, above=3pt] at ($(c1)!0.5!(c2)$) {$+U \sim N(3,1)$};
\node[efflbl] at ($(c1)!0.5!(c2) + (0, 1.3)$)
  {Cohen's $d = +1.871$\\(C2 vs C1)};

%% 분기 B: C1 → C3 (ㄷ자 직각 경로: 아래→오른쪽→위)
\draw[arr] (c1.south) -- ++(0,-0.6)
          -- ($(c3.south) + (0,-0.6)$) -- (c3.south);
\node[arrlbl, above=2pt] at ($(c1.south)!0.5!(c3.south) + (0,-0.6)$)
  {$+U \sim 6\!\cdot\!\mathrm{Pa}(3)$};
\node[efflbl] at ($(c2)!0.5!(c3) + (0, 1.3)$)
  {CVaR$_{99}$ 증가\\(C3 vs C2)};

% 하단 구성 요소
\node[comp, below=5mm of c0] {관측성: OFF\\불확실성: OFF};
\node[comp, below=5mm of c1] {관측성: ON\\불확실성: OFF};
\node[comp, below=5mm of c2] {관측성: ON\\불확실성: 가우시안(경꼬리)};
\node[comp, below=5mm of c3] {관측성: ON\\불확실성: 파레토(중꼬리)};

\end{tikzpicture}%
}% end resizebox
\caption{4개 실험 조건의 점진적 활성화}
\label{fig:conditions-progression}
\end{figure}


%% -------------------------------------------------------------------
\subsection{\chg{정적 기준선 C0}}
\label{subsec:c0}
%% -------------------------------------------------------------------

C0 조건은 관측성과 불확실성이 모두 비활성화된 상태로서,
기존 \nrev{ISO 22301:2019} 기반 정적 RTO 체계를 재현한다.
$O_{\text{raw}} = 0$, $U_{\text{raw}} = 0$이므로 $z_O = z_U = 0$이 되어
$S(0) = 0$이 성립하며, 따라서 $\DRTO = R_0 = 6.0$h, $\eta = 1.000$이다.
C0는 후속 조건들의 비교 기준선(baseline)으로 기능하며,
모든 동적 효과의 출발점을 제공한다.


%% -------------------------------------------------------------------
\subsection{\chg{C1의 관측성 기반 동적 조정}}
\label{subsec:c1}
%% -------------------------------------------------------------------

C1 조건은 관측성만 활성화하고 불확실성은 비활성화한 상태이다.
$U_{\text{raw}} = 0$이므로 $E_{U,\text{bnd}} = 0$이 되어
$\DRTO = R_0 - R_0 \cdot S(\kappa \cdot k_O \cdot O_{\text{raw}})$이다.
$S(\cdot) \geq 0$이므로 $\DRTO \leq R_0$가 보장되며,
이는 관측성 활성화가 복구 시간을 단축하는 방향으로만 작용함을 의미한다.
개별 바운딩의 핵심 장점으로, C1에서는 $R_{\max}$가 수식에 등장하지 않아
관측성만의 효과를 불확실성 상한과 독립적으로 분석할 수 있다.

$k_O$와 $O_{\text{raw}}$가 모두 최대값(1.0)일 때
$z_O = 1.2$가 되어
$S(1.2) \approx 0.537$이므로
$\DRTO \approx 2.78$h가 된다.
이는 완전 관측 상태에서 약 $53.7\%$의 복구 시간 단축이 가능함을 나타낸다.

\chg{이처럼 C1에서는 관측성 활성화가 복구 시간을 단축하는 방향으로만 작용하며, 그 평균적 단축 효과와 효과 크기가 기준선(C0) 대비 통계적으로 유의한가는 본 장이 제시한 이론적 예측에 해당한다. 표본 기반 실증값(평균 효율비, 단축률, Cohen's $d$)은 제4장에서 정량적으로 검증한다(\secref{sec:result_c0c1} 참조).}


%% -------------------------------------------------------------------
\subsection{\chg{C2의 가우시안 불확실성 결합}}
\label{subsec:c2}
%% -------------------------------------------------------------------

C2 조건은 관측성에 더하여 가우시안(Gaussian) 불확실성을 결합한 상태이다.
불확실성 원시값은 $U_{\text{raw}}^{(G)} \sim N(3, 1)$으로 생성되며,
$z_U = \kappa \cdot k_U \cdot R_0 \cdot U_{\text{raw}}^{(G)} / (R_{\max} - R_0)$으로 변환한 후
수식~(\ref{eq:E_U_bnd})에 따라 불확실성 효과 $E_{U,\text{bnd}}$를 산출한다.

\chg{한편 가우시안 분포는 음수 표본을 낼 수 있어 그 처리를 짚어 둔다.} 가우시안 분포의 지지(support)는 $(-\infty, +\infty)$이므로
$U_{\text{raw}}^{(G)} < 0$인 표본이 이론적으로 발생할 수 있다.
그러나 $N(3,1)$에서 $\Pr(U < 0) = \Phi(-3) \approx 0.135\%$로서
$n = 10{,}000$ 표본 중 13건에 불과하다.
해당 표본에서는 $z_U < 0$이 되어 $E_{U,\text{bnd}} < 0$이 산출되나,
실측 범위는 $[-1.13,\; -0.01]$h에 그치며
$\DRTO$ 최솟값이 $2.861$h로 물리적 하한($0$)을 충분히 상회한다.
따라서 절단정규분포(truncated normal)나 클리핑(clipping) 등의
인위적 분포 변형을 적용하지 않는다.
\chg{이는 세 가지 근거에 의한다. 첫째, 빈도 측면에서 음수 표본은 $0.13\%$(실측
13/10{,}000, 이론 $0.135\%$)에 불과하여 통계적 결론에 무시할 수 있는 영향만 미친다.
둘째, 경계 측면에서 해당 표본에서도 $\DRTO = 2.861$h로 0보다 충분히 크다. 셋째, 분포
보존 측면에서 절단을 적용하면 $N(3,1)$의 확률 구조가 왜곡되어 C2와 C3 사이의 공정한
비교(동일 기대값 $3.0$ 설계)가 훼손된다.}

가우시안 분포는 대칭적이고 경꼬리(light-tail) 특성을 가지므로,
C2 조건은 일상적 운영 변동, \chg{곧 인력 가용성의 일시적 변화와
장비 성능의 통상적 편차, 절차 수행 시간의 자연적 변이를} 반영한다.

\chg{C2에서는 이러한 가우시안 불확실성의 결합이 복구 시간을 연장하는 방향으로 작용하리라 예측되며, 그 평균적 연장 크기와 C1 대비 효과 크기는 제4장에서 정량적으로 검증한다(\secref{sec:result_c2} 참조).}


%% -------------------------------------------------------------------
\subsection{\chg{C3의 파레토 Heavy-tail 결합}}
\label{subsec:c3}
%% -------------------------------------------------------------------

C3 조건은 C1에서 분기하여 가우시안(C2) 대신 파레토(중꼬리, heavy-tail) 분포로
불확실성을 모델링한 변종이다. 불확실성 원시값은
$U_{\text{raw}}^{(P)} \sim 6 \cdot \mathrm{Pareto}(\alpha = 3)$으로
생성되며, $z_U = \kappa \cdot k_U \cdot R_0 \cdot U_{\text{raw}}^{(P)} / (R_{\max} - R_0)$으로 변환된다.
$\alpha = 3$은 평균은 유한하나 분산이 가우시안($\sigma^2=1$) 대비 약 27배인 분포를 제공한다.

C2와 C3의 불확실성 분포는 모집단 평균($\mu = 3.0$)이 동일하도록 설계되어
공정한 비교가 가능하다. 두 분포의 차이는 분산과 극단 백분위에서 발현된다.
가우시안의 99백분위가 약 5.33인 반면, 파레토의 99백분위는 약 21.85로서
약 4.1배의 차이를 보인다.

\chg{파레토 분포는 중꼬리 특성을 가지므로, C3는 평균 수준에서는 C2와 유사하더라도 극단 백분위로 갈수록 $\DRTO$가 C2를 크게 상회하리라 예측된다. 두 분포의 누적분포함수가 교차하는 지점이 존재하며, 본 연구에서는 이를 ``CDF 교차점(crossover)''으로 명명한다. 교차점을 경계로 핵심 영역(Core)과 꼬리 영역(Tail)을 분할하면, 꼬리 영역에서는 시그모이드 포화로 인해 관측성 가중치 $k_O$의 회귀 계수가 감쇠하는 이중채널 비선형 상호작용(dual-channel nonlinear interaction)이 나타나리라 예측된다. 교차점의 구체적 위치와 조건별 평균 효율비, 그리고 감쇠의 실증값은 \jrev{4.10절}(이중채널)에서 상세히 분석한다.}
